\chap{Введение}

Цели данной лабораторной работы:
\begin{enumerate}
	\item изучить назначение основных системных каталогов;
	\item приобрести навыки работы с базовыми командами для управления файлами и каталогами в операционных системах семейства GNU/Linux;
	\item научиться использовать встроенную справочную систему man.
\end{enumerate}

\chap{Выполнение}

\sect{Задание 1. Изучение структуры файловой системы}


\begin{figure}[!htbp]
	\centering
	%\includegraphics[width=0.6\linewidth]{""}
	\caption{}
\end{figure}
\FloatBarrier

%\lstinputlisting[style=python, caption={}]{}

\sect{Задание 2. Работа с файлами и каталогами}

\sect{Дополнительное задание 1}

\sect{Дополнительное задание 2. Анализ файлов системных каталогов}

\sect{Дополнительное задание 3. Сравнение содержимого системных каталогов}

\chap{Вывод}

В ходе лабораторной работы были: 
\begin{itemize}
	\item изучены назначения основных системных каталогов;
	\item приобретены навыки работы с базовыми командами для управления файлами и каталогами в операционных системах семейства GNU/Linux;
	\item полуяеяны знания об использовании встроенную справочную систему man.
\end{itemize}

\chap{Контрольные вопросы}

\begin{enumerate}
	\item \textit{В чем разница между операцией > и >> ? }\newline 
	\item \textit{Для чего используется каталог /home?}\newline 
	\item \textit{Какую информацию можно получить из каталога /proc?}\newline 
	\item \textit{Какой командой можно перейти из одного каталога в терминале в другой? Приведите пример команды перехода из каталога пользователя usr1 в корневой каталог}\newline 
	\item \textit{В чем разница между файлом и каталогом?}\newline 
	\item \textit{Как можно создать новый каталог?}\newline 
	\item \textit{Какая команда выдает справку по другой команде? }\newline 
	\item \textit{Как можно получить документацию по системному вызову exit? По команде exit?}\newline 
	\item \textit{В чем разница между жесткими и символическими ссылками?}\newline 
	\item \textit{Что произойдет с жесткой ссылкой, созданной командой: ln file1 link1, если удалить файл file1: rm file1?}\newline 
\end{enumerate}

